% LUT core-module ablation table, v1 (documentation only; no experiments were run to build it).
% Build: latexmk -pdf -interaction=nonstopmode -halt-on-error lut_ablation_table_v1.tex
\documentclass[10pt]{article}
\usepackage[landscape,margin=0.55in]{geometry}
\usepackage{lmodern}
\usepackage[T1]{fontenc}
\usepackage{amsmath,amssymb}
\usepackage{booktabs,longtable,array}
\usepackage{xcolor}
\usepackage[hang,flushmargin]{footmisc}
\usepackage{hyperref}
\hypersetup{colorlinks=true,linkcolor=black,urlcolor=blue}

\newcommand{\tbd}[1]{\textcolor{red!70!black}{\textbf{TBD}~--- #1}}
% \code{...}: typewriter, with a line-break opportunity after every underscore (long class names)
\newcommand{\code}[1]{{\renewcommand{\_}{\textunderscore\allowbreak}\texttt{#1}}}
\newcommand{\sgn}{\operatorname{sgn}}
\renewcommand{\arraystretch}{1.25}
\setlength{\LTpre}{4pt}\setlength{\LTpost}{4pt}

\title{LUT core-module ablation: three generations of lookup math\\[2pt]
\large Comparison table, version 1}
\date{2026-09-11}
\author{}

\begin{document}
\maketitle
\vspace{-2.5em}

\section*{Setup and notation}
\textbf{Shared model.} LUT transformer, $E=384$, 6 layers, 6 attention heads. The FFN is
\code{CompressionMultiHeadLUT} with $H=8$ heads, $tph=64$ tables per head, $d_{in}=d_{out}=48$:
\code{compress} $E\to H d_{in}$, the \emph{core LUT} on $[T,H,d_{in}]\to[T,H,d_{out}]$, \code{decompress}
$H d_{out}\to E$. Only the core LUT differs between rows. Per head $h$, the core reads $z_h\in\mathbb R^{d_{in}}$.

\textbf{Per table} $t$ (of the $tph$ tables of head $h$): $n$ fixed anchor pairs $(a_j,b_j)\subset[0,d_{in})$,
$u_j=z_h[a_j]-z_h[b_j]$, bits $\beta_j=[u_j>0]$, cell $c=\mathrm{pack}(\beta)\in[0,K)$ with $K=2^n$,
table $W_t\in\mathbb R^{K\times d_{out}}$. $c^{(j)}$ is $c$ with bit $j$ flipped, and
$j^\star=\arg\min_j|u_j|$ is the least-confident pair. The head output is $y_h=\sum_{t\in h}(\cdot)$.
In the backward, $g=\partial L/\partial y_h$, $G_t(k)=\langle g,W_t[k]\rangle$, and any
$\partial L/\partial u_j$ is scattered as $+$ to $z_h[a_j]$ and $-$ to $z_h[b_j]$. $\sigma$ is the logistic sigmoid.
The cell index is integer-valued: no gradient ever flows through it (no straight-through estimator in any row).

\textbf{Alternatives $k$.} Throughout, $k$ counts the table cells whose rows enter the computation, \emph{including}
the addressed cell: ``two alternatives'' is $k=2$ ($c$ and one flip), ``all alternatives'' is $k=K$.
\textbf{The cost also depends on $n$} (anchor pairs per table), which the variable list in the brief does not include.
All concrete numbers use $n=8$ ($K=256$), the value in the existing Gen-3 runs, and $P=H\cdot tph=512$ tables per token.

\section*{Counting convention for FLOPs and MACs}
\begin{itemize}\setlength\itemsep{1pt}
\item \textbf{Scope:} the core LUT module only, \textbf{per token}, one layer. Multiply by $T$ for a sequence and by 6 for the
  model. The \code{compress}/\code{decompress} projections are identical in every row and excluded: they add
  $E\cdot Hd_{in}+Hd_{out}\cdot E=294{,}912$ MACs ($\approx$\,589,824 FLOPs) forward and about twice that backward per token per
  layer, which is \emph{more} than the forward of every core variant below.
\item \textbf{FLOP} = one scalar floating-point operation: $+,-,\times,\div$, a float comparison, $|\cdot|$, $\exp$, $\log\sigma$,
  $\sigma$ each count 1. \textbf{MAC} = a multiply whose result is accumulated (weight $\times$ value, then added):
  \textbf{1 MAC = 2 FLOPs}, and MACs are included in the FLOP column. A pure accumulate (adding a looked-up row) is 1 FLOP and not a MAC.
\item \textbf{Not counted:} memory reads (table-row gathers, anchor gathers), integer work (bit packing, XOR, index offsets),
  optimizer updates, and saving activations. Row reads per token are $P\cdot k$ rows of $d_{out}$ values where they matter.
\item \textbf{Forward} = inference cost (what an evaluated model executes). \textbf{Backward} = all \emph{additional}
  arithmetic a training step needs, including training-only forward work (for example the alternative search in 1.1).
\item Terms proportional to $d_{out}$, $K$ and $nK$ are exact. Scalar per-table overheads ($O(n)$ and constants) are
  counted from the code but approximate to within a few operations per table, since autograd's internal breakdown varies.
  $d_{in}$ does not enter the core cost: it only sets the anchor index range.
\item Temperature gradients ($O(1)$ per layer) are ignored. Gen-2 rows assume the confidence gate is \emph{off}, as in the existing Gen-2 runs.
\end{itemize}

\newpage
\section*{Table A --- forward and backward math}
\footnotesize
\begin{longtable}{>{\raggedright\arraybackslash}p{2.9cm} >{\raggedright\arraybackslash}p{8.8cm} >{\raggedright\arraybackslash}p{12.0cm}}
\toprule
\textbf{Variant} & \textbf{Forward pass (math)} & \textbf{Backward pass (math)} \\
\midrule
\endfirsthead
\toprule
\textbf{Variant} & \textbf{Forward pass (math)} & \textbf{Backward pass (math)} \\
\midrule
\endhead
\midrule \multicolumn{3}{r}{\emph{continued on next page}} \\
\endfoot
\bottomrule
\endlastfoot

\multicolumn{3}{l}{\textbf{GEN 1 --- Spiking Manifesto basic model, rational uncertainty $U(u)=0.5/(1+|u|)$}\footnotemark[1]} \\[2pt]
1.1 Hard forward, two-alternatives backward ($k=2$) &
$y_h=\sum_t W_t[c_t]$. \newline
Training only: $j^\star=\arg\min_j|u_j|$ (for $k>2$: the $k-1$ smallest $|u_j|$). &
Weights: $\partial L/\partial W_t[c_t]\mathrel{+}=g$ (addressed row only). \newline
Input (surrogate --- the derivative of 1.2's blend): $\partial L/\partial u_{j^\star}=-U'(u_{j^\star})\,\big(G_t(c_t)-G_t(c_t^{(j^\star)})\big)$,
with $-U'(u)=\dfrac{0.5\,\sgn u}{(1+|u|)^2}$; $\partial L/\partial u_j=0$ for $j\neq j^\star$. For $k>2$ each flip contributes and the sum is divided by $k-1$. \\[3pt]
1.2 Soft forward, two-alternatives backward ($k=2$) &
$y_h=\sum_t\big[(1-U_t)\,W_t[c_t]+U_t\,W_t[c_t^{(j^\star)}]\big]$, $U_t=U(u_{t,j^\star})$. \newline
Continuous at $u_{j^\star}=0$ (both weights $\tfrac12$). Applied at eval too. &
Weights: $\partial L/\partial W_t[c_t]\mathrel{+}=(1-U_t)\,g$, $\partial L/\partial W_t[c_t^{(j^\star)}]\mathrel{+}=U_t\,g$. \newline
Input: same expression as 1.1, which here is the \emph{exact} gradient of the blend with the two cells held fixed. \\[4pt]
\midrule
\multicolumn{3}{l}{\textbf{GEN 2 --- FastMHL math}. Surrogate: $\rho_j=\dfrac{|u_j|}{T_s+|u_j|}$,
$s_k=\sum_j\rho_j-2\!\!\sum_{j:\,\mathrm{bit}_j(k)\neq\beta_j}\!\!\rho_j$, $\pi_k=\mathrm{softmax}_k(s_k/T_{sel})$ (signs $\beta$ pinned; $T_s,T_{sel}$ are temperatures)} \\[2pt]
2.1 Hard forward, all-alternatives backward ($k=K$) &
$y_h=\sum_t W_t[c_t]$. \newline
($c_t=\arg\max_k s_k$; no $K$-sized tensor is built.) &
Weights: $\partial L/\partial W_t[c_t]\mathrel{+}=g$. \newline
Input via the $K$-cell surrogate $\hat y_t=\sum_k\pi_kW_t[k]$ (never used in the forward):
$\dfrac{\partial L}{\partial u_j}=\dfrac{1}{T_{sel}}\sum_{k}\pi_k\big(G_t(k)-\bar G_t\big)\dfrac{\partial s_k}{\partial\rho_j}\cdot\dfrac{T_s\,\sgn u_j}{(T_s+|u_j|)^2}$,
$\bar G_t=\sum_k\pi_kG_t(k)$, $\partial s_k/\partial\rho_j=\pm1$ (agree/disagree with $\beta_j$). All $n$ pairs get gradient; reads all $K$ rows. \\[3pt]
2.2 Hybrid\_smooth two-alternatives forward, all-alternatives backward (fwd $k=2$, bwd $k=K$) &
$y_h=\sum_t\big[(1-w_t)W_t[c_t]+w_tW_t[c_t^{(j^\star)}]\big]$, \newline
$w_t=\sigma(-\Delta_t/T_{sel})$, $\Delta_t=2\rho_{t,j^\star}$ \newline
(the exact 2-way softmax of $s$ over $\{c,c^{(j^\star)}\}$). Applied at eval too. &
Weights: $(1-w_t)\,g$ into $W_t[c_t]$ and $w_t\,g$ into $W_t[c_t^{(j^\star)}]$. \newline
Input: the $K$-cell surrogate of 2.1 (not the gradient of the forward that actually ran). \\[3pt]
2.3 Hard forward, 2-alternatives backward &
$y_h=\sum_t W_t[c_t]$ &
\tbd{not implemented.} \code{FastMultiHeadLut} raises \code{ValueError} for \code{backward\_topk>0} unless
\code{forward\_mode='hybrid\_smooth'} (\code{fast\_multi\_head\_lut.py:1305}). \\[3pt]
2.4 Hybrid\_smooth two-alternatives forward, 2-alternatives backward ($k=2$) &
As 2.2. &
Weights: as 2.2. \newline
Input (\code{backward\_topk=1}): the exact gradient of the 2-cell blend with cells fixed. Only $j^\star$ gets gradient:
$\dfrac{\partial L}{\partial u_{j^\star}}=-\dfrac{w_t(1-w_t)}{T_{sel}}\big(G_t(c_t^{(j^\star)})-G_t(c_t)\big)\dfrac{2T_s\,\sgn u_{j^\star}}{(T_s+|u_{j^\star}|)^2}$. \\[4pt]
\midrule
\multicolumn{3}{l}{\textbf{GEN 3 --- LightMHL math (LookupFFN-inspired)}. Margin score
$\,s_t=\big(\sum_j|u_{t,j}|\big)\prod_j\sigma(2|u_{t,j}|)$; the cell index uses detached $u$\footnotemark[2]} \\[2pt]
3.1 Soft forward $n{=}1$, margin confidence, single-alternative backward ($k=1$) &
$y_h=\sum_t s_t\,W_t[c_t]$. \newline
One hard cell read scaled by a differentiable score: $y$ jumps at every cell boundary. &
Weights: $\partial L/\partial W_t[c_t]\mathrel{+}=s_t\,g$. \newline
Input (only through the score): $\dfrac{\partial L}{\partial u_j}=G_t(c_t)\Big(\prod_i\sigma(2|u_i|)+2s_t\,\sigma(-2|u_j|)\Big)\sgn u_j$
for all $j$. No other cell is compared. \\[3pt]
3.2 Soft forward $n{=}2$, margin confidence, two-alternatives backward ($k=2$) &
$y_h=\sum_t s_t\big[\omega_0W_t[c_t]+\omega_1W_t[c_t^{(j^\star)}]\big]$, \newline
$\omega_1=\sigma(-2|u_{j^\star}|/\tau)$, $\omega_0=1-\omega_1$, $\tau=e^{\log\tau}$ (learnable). &
Weights: $s_t\omega_0\,g$ and $s_t\omega_1\,g$ into the two read cells. \newline
Input: $\dfrac{\partial L}{\partial u_j}=(\omega_0G_0+\omega_1G_1)\dfrac{\partial s_t}{\partial|u_j|}\sgn u_j+[j=j^\star]\,s_t\tfrac{2}{\tau}\omega_0\omega_1(G_0-G_1)\sgn u_j$,
with $G_0=G_t(c_t)$, $G_1=G_t(c_t^{(j^\star)})$. \newline
$\partial L/\partial\log\tau=\sum_t s_t\omega_0\omega_1(G_1-G_0)\,2|u_{j^\star}|/\tau$. \\
\end{longtable}
\normalsize
\footnotetext[1]{Other uncertainty functions may be tried later. The code also offers \code{UncertaintyMode.INVERSE\_QUADRATIC}, $0.5/(1+u^2)$.}
\footnotetext[2]{A hard-forward Gen-3 variant still needs to be designed: nothing in the code trains LightMHL with a hard forward.}

\newpage
\section*{Table B --- cost and results}
Per token, one layer, core LUT only. Concrete values at $H=8$, $tph=64$, $d_{in}=d_{out}=48$, $n=8$ ($K=256$), $P=H\,tph=512$, $d=d_{out}$.
Validation bpb is on 16K steps, batch 48 rows $\times$ 512 tokens.
\footnotesize
\setlength{\tabcolsep}{3pt}% seven columns: narrower gutters keep the table inside the text block
\begin{longtable}{>{\raggedright\arraybackslash}p{1.3cm} >{\raggedright\arraybackslash}p{2.8cm} >{\raggedright\arraybackslash}p{4.2cm} >{\raggedright\arraybackslash}p{2.0cm} >{\raggedright\arraybackslash}p{2.8cm} >{\raggedright\arraybackslash}p{4.3cm} >{\raggedright\arraybackslash}p{6.0cm}}
\toprule
\textbf{Variant} & \textbf{Forward FLOPs} & \textbf{Backward FLOPs} & \textbf{Forward MACs} & \textbf{Backward MACs} & \textbf{Validation bpb (16K, bs48)} & \textbf{Implementation (class used)} \\
\midrule
\endfirsthead
\toprule
\textbf{Variant} & \textbf{Forward FLOPs} & \textbf{Backward FLOPs} & \textbf{Forward MACs} & \textbf{Backward MACs} & \textbf{Validation bpb (16K, bs48)} & \textbf{Implementation (class used)} \\
\midrule
\endhead
\midrule \multicolumn{7}{r}{\emph{continued on next page}} \\
\endfoot
\bottomrule
\endlastfoot
1.1 &
$P(2n+d)$ \newline $=\mathbf{32{,}768}$ &
$P\big(2kd+d+2n+9(k{-}1)\big)$ \newline $=\mathbf{135{,}680}$ &
$0$ &
$P\,kd$ \newline $=\mathbf{49{,}152}$ &
\tbd{not yet implemented} in the FFN harness; no run exists. &
\code{spiky.lutorch.multi\_head\_lut.MultiHeadLut} (\code{smooth\_mode=False}, \code{n\_alternatives=1}, \code{uncertainty\_mode=INVERSE\_L1}); also \code{lut\_fused.LUTLayer.LUTLayerBasic}. \textbf{Not wired into} \code{CompressionMultiHeadLUT}. \\
1.2 &
$P(2kd+4n+3k-2)$ \newline $=\mathbf{116{,}736}$ &
$P\big(4kd+9(k{-}1)\big)$ \newline $=\mathbf{201{,}216}$ &
$P\,kd$ \newline $=\mathbf{49{,}152}$ &
$P\,2kd$ \newline $=\mathbf{98{,}304}$ &
\tbd{not yet implemented} in the FFN harness; no run exists. &
\code{MultiHeadLut} (\code{smooth\_mode=True}, \code{n\_alternatives=1}). Not wired into \code{CompressionMultiHeadLUT}. \\
\midrule
2.1 &
$P(2n+d)$ \newline $=\mathbf{32{,}768}$ &
$P(2Kd+4nK+9K+11n+d)$ \newline $=\mathbf{18{,}026{,}496}$\footnotemark[3] &
$0$ &
$P(Kd+2nK)$ \newline $=\mathbf{8{,}388{,}608}$ &
\tbd{not yet run} on the paper protocol. \newline
Prior: \textbf{1.2281} (\code{exp\_g\_0014}), $n{=}6$, old eval\footnotemark[4] &
\code{CompressionMultiHeadLUT(lut\_impl='fast')} $\to$ \code{FastMultiHeadLut(forward\_mode='hard', backward\_topk=0)}; keys \code{lut\_forward\_mode}, \code{lut\_backward\_topk}. \\
2.2 &
$P(2kd+4n+5)$ \newline $=\mathbf{117{,}248}$ &
$P(2Kd+4nK+9K+11n+2kd)$ \newline $=\mathbf{18{,}100{,}224}$\footnotemark[3] &
$P\,kd$ \newline $=\mathbf{49{,}152}$ &
$P(Kd+2nK+kd)$ \newline $=\mathbf{8{,}437{,}760}$ &
\tbd{not yet run} on the paper protocol. \newline
Prior: \textbf{1.2176} (\code{exp\_n\_0044}), $n{=}6$, old eval\footnotemark[4] &
\code{FastMultiHeadLut(forward\_mode='hybrid\_smooth', backward\_topk=0)}. \\
2.3 &
$P(2n+d)=32{,}768$ &
n/a &
$0$ &
n/a &
\tbd{not yet implemented} &
Would be \code{FastMultiHeadLut(forward\_mode='hard', backward\_topk=1)}, which the code rejects. \\
2.4 &
$P(2kd+4n+5)$ \newline $=\mathbf{117{,}248}$ &
$P(4kd+9n+15)$ \newline $=\mathbf{241{,}152}$ &
$P\,kd$ \newline $=\mathbf{49{,}152}$ &
$P\,2kd$ \newline $=\mathbf{98{,}304}$ &
\tbd{not yet run} (the only \code{backward\_topk>0} run, \code{exp\_n\_0188}, is $H{=}4$, $tph{=}128$, $n{=}7$). &
\code{FastMultiHeadLut(forward\_mode='hybrid\_smooth', backward\_topk=1)}. Note: \code{backward\_topk} counts cells \emph{beyond} the addressed one. \\
\midrule
3.1 &
$P(2kd+7n)$ \newline $=\mathbf{77{,}824}$ &
$P(4kd+8n)$ \newline $=\mathbf{131{,}072}$ &
$P\,kd$ \newline $=\mathbf{24{,}576}$ &
$P\,2kd$ \newline $=\mathbf{49{,}152}$ &
\tbd{not yet run} at 16K. \newline
48K: \textbf{1.1313} (\code{exp\_g\_0207})\footnotemark[5] &
\code{CompressionMultiHeadLUT(lut\_impl='light')} $\to$ \code{LightMultiHeadLUT(confidence\_form='margin', read\_top\_n=1)}. \\
3.2 &
$P(2kd+8n+7)$ \newline $=\mathbf{134{,}656}$ &
$P(4kd+8n+11)$ \newline $=\mathbf{235{,}008}$ &
$P\,kd$ \newline $=\mathbf{49{,}152}$ &
$P\,2kd$ \newline $=\mathbf{98{,}304}$ &
\tbd{not yet run} at 16K. \newline
48K: \textbf{1.1304} (\code{exp\_g\_0206})\footnotemark[5] &
\code{LightMultiHeadLUT(confidence\_form='margin', read\_top\_n=2, read\_tau=0.5, learnable)}. \\
\end{longtable}
\normalsize
\footnotetext[3]{Dominated by $2Kd$: the all-alternatives backward reads all $K$ rows of every table. At the existing runs' $n=6$ ($K=64$) the 2.1 backward is $P\cdot 8{,}370=4{,}285{,}440$ FLOPs.
Memory: at $n=8$ the backward holds several \code{[tokens, 512, 256]} fp32 buffers, about 12.9\,GiB each at 24,576 tokens, so batch 48 needs micro-batching on a 32\,GB card.}
\footnotetext[4]{Old batch-coupled eval (245,760 val tokens from token 0, no row skip), $n=6$, tied embeddings, learnable temperatures; branch \code{research/hyperplane\_ffn\_next}. \textbf{Not comparable} to the corrected protocol. Same-protocol neighbours: \code{exp\_n\_0052} 1.2286 (hard, batched temperatures).}
\footnotetext[5]{48,000-step run, $n=8$, untied, fixed eval; device batch 12 with gradient accumulation to 24,576 tokens. \textbf{Not a 16K result.} The step-16,000 rows (1.1862 for n=1, 1.1793 for n=2) sit on a 48K cosine schedule and are not comparable either.}

\section*{Where the code disagrees with the brief}
\begin{enumerate}\setlength\itemsep{1pt}
\item \textbf{Gen 1 is not wired into the experiment harness.} \code{CompressionMultiHeadLUT} accepts only
  \code{lut\_impl} $\in\{$\code{fast}, \code{light}, \code{bh4}$\}$ (\code{compression\_mhl.py:184}), and no \code{model\_build.py} key reaches
  \code{MultiHeadLut}. \code{MultiHeadLut} also takes a shared \code{[B, input\_dim]} input, not the per-head \code{[T,H,d\_in]} layout, so
  rows 1.1/1.2 need new wiring (for example \code{anchor\_candidates} restricting each head's anchors) before they can run.
\item \textbf{In 1.1 the input gradient is a surrogate.} It is the derivative of 1.2's soft blend, while the weights update only the addressed row.
\item \textbf{The paper's status report writes $u(\delta)=0.5/(T+|\delta|)$} with a per-LUT temperature. The code's
  $U(u)=0.5/(1+|u|)$ has no temperature (\code{l\_projection.py:96}, \code{lutorch.cu:797}).
\item \textbf{2.3 cannot be built}: hard forward with a sparse backward raises \code{ValueError} (\code{fast\_multi\_head\_lut.py:1305}).
\item \textbf{\code{backward\_topk} is off by one relative to ``$k$ alternatives''.} It counts cells beyond the addressed one:
  ``2 alternatives'' is \code{backward\_topk=1}, and \code{backward\_topk=2} means 3 cells.
\item \textbf{``All alternatives'' in Gen 2 is all $K=2^n$ cells} (a full softmax surrogate), not the $n$ single-bit flips (that would be \code{backward\_topk=n}).
  \code{backward\_topk} never changes the weight gradient, which always uses the rows the forward read.
\item \textbf{Gen 2's hybrid\_smooth weights are $w=\sigma(-2\rho_{j^\star}/T_{sel})$}, a temperature-controlled rational soft-sign, not Gen 1's $U(u)$. It does match ``two-alternatives forward''.
\item \textbf{3.1 ``soft forward $n{=}1$'' is not soft in the cell sense.} It reads one hard cell and scales it by a differentiable margin score, and its backward compares no alternative cell.
  ``Single-alternative backward'' is accurate only in the sense that the input gradient flows through the score alone.
\item \textbf{The margin confidence function is LookupFFN's kernel} $s=(\sum_j|u_j|)\prod_j\sigma(2|u_j|)$, not $\min_j|u_j|$. The factor 2 is fixed.
\item \textbf{The cost depends on $n$ (anchor pairs), which the brief's variable list omits.} Existing Gen-2 16K runs use $n=6$, Gen-3 runs use $n=8$, so a paper set needs one $n$ for all rows.
\item \textbf{Minor:} \code{exp\_g\_0207} (3.1) sets \code{lut\_read\_tau\_learnable=true}, which creates a parameter that never receives a gradient at $n=1$.
\end{enumerate}

\section*{TBD cells and experiments still to run}
\begin{itemize}\setlength\itemsep{1pt}
\item \textbf{Needs implementation first:} 1.1 and 1.2 (wire \code{MultiHeadLut} into \code{CompressionMultiHeadLUT}); 2.3 (a hard-forward path for \code{backward\_topk}).
\item \textbf{Needs a run on the paper protocol} (16K steps, batch 48, corrected eval, one fixed $n$): 2.1, 2.2, 2.4, 3.1, 3.2. The 2.1/2.2 numbers above are $n=6$ on the old eval, and 3.1/3.2 exist only at 48K.
\item \textbf{Every row's validation bpb is TBD on the paper protocol.} Cost cells for 2.3 are n/a until it exists.
\item \textbf{Reference anchors} (vanilla dense FFN, same protocol, corrected eval): 1.1651 / 1.1618 at 16K (two seeds), 1.1154 at 48K.
\end{itemize}
\end{document}
